\section{The full distinct-entry diagonal square modulo an extending pair}
\label{sec:r10-commuting-diagonal-relative-square}

\paragraph{Draft status.}
This module is submitted for independent exact-source review.  It treats
precisely the prescribed antisymmetric SHT primitives, with all their
metric terms and absolute-value walls.  Its inputs are R10 MT and the
frozen R9 PPC, CPT, CDP, RAG and RSP interfaces.

Let $N$ be an ordinary analytic open over the original $K$, let $L$
carry its fixed piecewise smooth metric with curvature $R$, and let
$f=(f_1,\ldots,f_r)$ be analytic sections of $L$.  No regularity is
needed for the local identity.  Put
\[
 C=|f=0|,\qquad W=N\setminus C,\qquad s=r-1,\qquad d=r-2.
 \tag{CS1}
\]
Let $\lambda,\mu$ be analytic units on all of $N$, write
$a=v(\lambda)$, $b=v(\mu)$, and let $A$ and $B$ multiply the first
and second entries respectively.  Write $T=AB=BA$.  The four
potentials $\Phi_0,\Phi_A,\Phi_B,\Phi_T$ and their projective
curvatures $C_i=R-D\Phi_i$ are actual on $W$.  Define
\[
 P_s=p_{0A,a}+p_{AT,b}-p_{0B,b}-p_{BT,a},\qquad
 e_{a,b}=(a\,d''b,\ b\,d'a).
 \tag{CS2}
\]
Every $p$ in (CS2) is exactly R9 SHT4--SHT6.  Equivalently use the
oriented loop $0\to A\to T\to B\to0$ with weights $a,b,-a,-b$.
The pair $e_{a,b}$ is actual and closed on all of $N$ by
$Da=Db=0$ and the sign cancellation in R9 CPT11.

We construct an actual triple $q_s$ on all of $W$ such that
\[
                    P_s=\partial_2q_s+e_{a,b}R^{s-1}|_W.
 \tag{CS3}
\]
For $r=2$ this is the frozen CDP theorem.  The new construction is
for every $r\ge3$, with arbitrary additional base directions.  The
last pair in (CS3) is defined across the common zero set.  We do not
require or assert a strict witness with that pair removed.

\subsection{A universal edge reduction with its full residual}

For an even diagonal form $Q$ and actual scalars $u,v$, write
\[
 \begin{split}
 s'(u,v)&=u\,d'v-v\,d'u,\qquad
 s''(u,v)=u\,d''v-v\,d''u,\\
 \mathfrak a(u,v;Q)&=(s''(u,v)Q,-s'(u,v)Q).
 \end{split}
 \tag{CS4}
\]
The second-component sign is part of the definition.  A formal
letter $d$ below denotes a coefficient calculation performed
separately in each colour; it never identifies $d'$ and $d''$.

Consider one rank-one unit edge with actual potentials $p,q$, unit
weight $c$, and a closed reference curvature $R_0$.  Retain
\[
 C_p=R_0-Dp,\qquad C_q=R_0-Dq,\qquad t=|c|,
 \tag{CS5}
\]
and the actual norm relation $p=q$ on $c=0$.  Set
\[
 \begin{gathered}
 H_k=h_{k-1}(R_0,C_p,C_q),\qquad J_k=h_{k-1}(C_p,C_q),\\
 K_{p,k}=h_{k-1}(R_0,C_p),\qquad K_{q,k}=h_{k-1}(R_0,C_q),\\
 M_k=-\frac14\bigl[(q^2-p^2)H_k+t(q-p)J_k
                                      +c(pK_{p,k}+qK_{q,k})\bigr].
 \end{gathered}
 \tag{CS6}
\]
Here $h_j$ is the complete homogeneous polynomial in all displayed
variables, with $h_0=1$ and $h_j=0$ for $j<0$.  Define
\[
 S_{e,k}=\mathfrak a(p,q;H_k)
             +\mathfrak a(c_+,q;K_{q,k})
             +\mathfrak a(c_-,p;K_{p,k}),\qquad
 c_+=\max(c,0),\quad c_-=\min(c,0).
 \tag{CS7}
\]
This pair has degrees $(k-1,k)$ and $(k,k-1)$.  We make no
separate closedness claim about it or its summands.

Let $Q_k^{\mathrm{tet}}(R_0;t,p,q)$ be the explicit PPC8--PPC10
triple for the four potentials $(0,t,p,q)$ in degree $k$.
For $k\ge2$, put
\[
 F_k=h_{k-2}(R_0,R_0-Dt,C_p,C_q)
 \tag{CS8}
\]
and define
\[
 \begin{split}
 A_k={}&(t\,d''p\wedge d''q-p\,d''t\wedge d''q
                                      +q\,d''t\wedge d''p)F_k,\\
 B_k={}&(t\,d'p\wedge d'q-p\,d't\wedge d'q
                                      +q\,d't\wedge d'p)F_k,\\
 Z_k={}&\bigl[t(d'p\wedge d''q-d'q\wedge d''p)
           -p(d't\wedge d''q-d'q\wedge d''t)\\
 &\hspace{24mm}+q(d't\wedge d''p-d'p\wedge d''t)\bigr]F_k,\\
 Q_k^{\mathrm{tet}}&=(A_k/3,B_k/3,-Z_k/6).
 \end{split}
 \tag{CS9}
\]
Set $Q_1^{\mathrm{tet}}=Q_0^{\mathrm{tet}}=0$.  Its boundary is
$\delta T_k=T_{t,p,q}-T_{0,p,q}+T_{0,t,q}-T_{0,t,p}$.
This definition retains the full curvature $R_0-D|c|$.

\begin{lemma}[Full edge formula]
For the literal SHT expression with degree $k$, including when it is
only a truncated expression rather than a rank-$r$ primitive,
\[
 \boxed{\quad
 p_e^{[k]}=
 \partial_2\left\{(0,0,M_k)+\frac12
       \bigl(Q_k^{\mathrm{tet}}-R_0Q_{k-1}^{\mathrm{tet}}\bigr)\right\}
                                      -\frac12S_{e,k}.
 \quad}
 \tag{CS10}
\]
All absolute-value walls remain.  No curvature coefficient pair is
declared closed.
\end{lemma}

\begin{proof}
All curvature factors in (CS6) are closed.  Flattening the two
first derivatives only for the coefficient calculation, the sum
of the SHT components is $2dM_k$ and their difference is
\[
 \theta_e=-s(p,q)H_k-\frac12s(t,q-p)J_k
               -\frac12\{s(c,p)K_{p,k}+s(c,q)K_{q,k}\},
 \qquad s(u,v)=u\,dv-v\,du.
 \tag{CS11}
\]
This is an expansion of the entire SHT6 expression: $H_k$ has the
reference curvature as an additional variable, so the lower metric
terms are included.

In the coefficient calculation for $\delta T_k$, one may replace
$R_0-Dt$ by $R_0$ in its curvature polynomials, for the following
actual reason.  Every term of their difference has a factor $Dt$,
whose pure polyhedral support lies on $c=0$.  On each carrier there,
$t=0$ with zero tangential first jet.  Also $q-p=0$ with zero
tangential first jet.  The symbol in $T_{t,p,q}$ consequently
vanishes because it is built from $q-p$; those in $T_{0,t,p}$ and
$T_{0,t,q}$ vanish because they contain $t$ and its first derivative.
Every term with a positive power of $Dt$ is killed by restricting
its scalar and first-form coefficient to its carrier.  The same
calculation holds on each further actual refinement.  It drops no
derivative-of-Dirac term.

The resulting difference of flattened components is
\[
 \delta\theta_k=s(t,p)\{H_k-h_{k-1}(R_0,R_0,C_p)\}
             -s(t,q)\{H_k-h_{k-1}(R_0,R_0,C_q)\}.
 \tag{CS12}
\]
The complete homogeneous identities
\[
 \begin{split}
 H_k-R_0H_{k-1}&=J_k,\\
 h_{k-1}(R_0,R_0,C_p)
     -R_0h_{k-2}(R_0,R_0,C_p)&=K_{p,k}
 \end{split}
 \tag{CS13}
\]
and the analogous $q$ identity give
\[
 \delta\theta_k-R_0\delta\theta_{k-1}
       =s(t,p)(J_k-K_{p,k})-s(t,q)(J_k-K_{q,k}).
 \tag{CS14}
\]
For $k=1$ both tetrahedron differences are zero and this reads
$0=0$.  Subtract one half of (CS14) from (CS11), and use
$(t+c)/2=c_+$ and $(c-t)/2=c_-$.  The remainder is exactly
the negative of the first-colour expression defining $S_{e,k}$.
The other colour has its prescribed opposite sign.  Reconstituting
the pair gives the scalar boundary, one half of the tetrahedron
pair boundary, and $-S_{e,k}/2$, proving (CS10).
\end{proof}

\subsection{The two support identities for the residual}

The following arguments concern actual projective curvatures.  On
a local unit presentation the tuple norm is a finite minimum of
valuation coordinates with the common metric added.  Its projective
curvature is the pure polyhedral corner of that minimum.  It
vanishes on the ordinary open where one entry is strictly the unique
minimizer.  In curvature products the scalar and first-form
coefficients are restricted tangentially to each carrier.  Equality
of scalars on an entire carrier therefore gives equality of their
full restricted first jets.  These are the frozen RAT3/TGC6 support
interfaces, with all actual refinement relations retained.

For an edge raising one entry by $c\ge0$, every point in the support
of the old corner has an unchanged entry attaining the old minimum.
It still attains that minimum after the raise, so
\[
                         p=q\quad\hbox{on }\operatorname{supp}C_p.
 \tag{CS15}
\]
This is unaffected by a common reference potential.  Thus every
$\mathfrak a(p,q;C_p^iQ)$ with $i\ge1$ is zero, with full tangential
first jets retained.  On a carrier where the changed entry attains
the new minimum, its old value is the old minimum, and hence
\[
                              p+c=q.
 \tag{CS16}
\]
For $c\le0$ apply the same facts to the reverse edge.  On the wall
$c=0$, one has $p=q$, while the restrictions of $c,c_+,c_-$ and
their tangential first derivatives are zero.  These statements
are checked on every carrier, including intersections with other
walls.  No sign-chamber formula is extended by a characteristic
function.

\begin{lemma}[Two-entry recurrence]
Suppose on an ordinary open all four square norms are determined
by the first two entries.  Let $R_0$ be any common closed reference
curvature, with the normalized potentials satisfying (CS5).
For every edge and $k\ge2$,
\[
                         S_{e,k}=R_0S_{e,k-1}.
 \tag{CS17}
\]
Consequently the full square satisfies
\[
                \sum_e S_{e,s}
                      =-\mathfrak a(a,b;R_0^{s-1}).
 \tag{CS18}
\]
\end{lemma}

\begin{proof}
The difference in (CS17), by (CS13), is
\[
 \mathfrak a(p,q;h_{k-1}(C_p,C_q))
       +\mathfrak a(c_+,q;C_q^{k-1})
       +\mathfrak a(c_-,p;C_p^{k-1}).
 \tag{CS19}
\]
For $c\ge0$, (CS15) kills every monomial containing $C_p$.
The remaining term is $\mathfrak a(p+c,q;C_q^{k-1})$.
Since $k-1\ge1$, its carrier lies in the new two-entry corner,
where both entries attain the new minimum.  Equation (CS16)
kills the symbol and its restricted first jet there.  Where one
entry vanishes and the other is active, the corner is already
zero on an ordinary neighborhood.  Reverse the edge for $c\le0$;
the preceding wall check covers every carrier in $c=0$.
This proves (CS17).

In degree one all four polynomials $H_1,J_1,K_{p,1},K_{q,1}$
are one and the tetrahedron term is zero.  The full rank-two
CDP8--CDP12 calculation, including zero-entry neighborhoods and an
arbitrary common scalar or metric, gives
\[
                          \sum_e S_{e,1}=-\mathfrak a(a,b;1).
 \tag{CS20}
\]
It can also be checked by substituting the four two-entry minima
in (CS7) and using the four homogeneous first-jet relations CDP10.
Iterating (CS17) proves (CS18).  Mixed products $C_pC_q$ have
not been declared zero: the coefficient relations kill their
specific symbol factors before the other curvatures are wedged.
\end{proof}

\begin{lemma}[Passive-block normalization]
Let $r\ge3$ and suppose the passive tuple $(f_3,\ldots,f_r)$
is nonzero on an ordinary open.  Normalize by its norm potential,
so its maximum norm is one and its projective curvature is the
reference $R_0$.  For $s=r-1$,
\[
                    R_0^{s-1}=0,\qquad S_{e,s}=0
                                \quad\hbox{for every square edge}.
 \tag{CS21}
\]
\end{lemma}

\begin{proof}
The reference curvature is the pullback of the standard projective
form by the genuine original-$K$ map defined by the $r-2$ passive
sections, whose target is $\mathbf P^{r-3}$.  Its $(r-2)$nd power
is zero, with all mixed base coefficients retained.  For $r=3$
the target is $\mathbf P^0$ and the reference curvature is zero.

For $c\ge0$, (CS15) removes every $C_p$-containing monomial from
the first term of (CS7).  What remains is
\[
             \mathfrak a(p+c,q;h_{s-1}(R_0,C_q)).
 \tag{CS22}
\]
Its pure $R_0^{s-1}$ term is zero.  Every other monomial contains
a positive power of $C_q$.  On a carrier of such a monomial,
if the changed entry is active, use (CS16).  If it is inactive
and a passive entry is active, the normalized potential is $q=0$
with zero tangential first jet.  If both the changed entry and
all passive entries are inactive, the other affected entry is
the unique minimum, so $C_q$ vanishes on an ordinary neighborhood
and this case contributes no carrier.  These cases exhaust the
active set.  Reverse the edge for $c\le0$ and retain the wall
check at zero.  This proves (CS21) as an actual pair identity.
\end{proof}

\subsection{One capped actual norm on the entire punctured overlap}

The passive norm potential diverges at passive-zero points, which
may lie in $W$.  A single capped potential avoids any extension
of triples across that locus.  Let
\[
 \psi=-\log\|(f_3,\ldots,f_r)\|_{L,\infty},
 \tag{CS23}
\]
interpreted as $+\infty$ at the passive zero set.  Choose one real
number $\varepsilon>0$ and put
\[
 \kappa=|a|+|b|+\varepsilon,\qquad
                 \Psi=\min(\psi,\Phi_0+\kappa)\quad\hbox{on }W.
 \tag{CS24}
\]
This is an ordinary actual scalar on all of $W$.  On the
passive-nonzero locus it has a finite piecewise smooth presentation.
Near a passive-zero point, some affected section is nonzero.  The
ratio of the passive maximum norm to the full maximum norm is a
continuous function vanishing there, while $e^{-\kappa}$ is
continuous and positive.  On an ordinary neighborhood that ratio
is less than $e^{-\kappa}$, so $\Psi=\Phi_0+\kappa$ literally.
This defines an actual scalar across the passive zero set by
equality on an open overlap.

The common metric term occurs once in each entry of the minimum
and cancels from their comparison.  On every fixed whole local
unit presentation, the cap is thus one finite piecewise affine
operation together with the original metric potential.  All
real offsets and absolute-value walls remain.  No element of $K$
with valuation $\varepsilon$ is required.  Define
\[
                 R_*=R-D\Psi,\qquad p_i=\Phi_i-\Psi.
 \tag{CS25}
\]
Then $C_i=R_*-Dp_i$ are the original projective curvatures.
Let $P_s^*$ be the full SHT loop for these potentials and reference
curvature.  MT16 with common scalar $h=-\Psi$ gives an actual
triple $Q_{\rm met}$ on all of $W$ satisfying
\[
                   P_s^*-P_s=\partial_2Q_{\rm met}.
 \tag{CS26}
\]
There is no requirement that $\Psi$ extend through the common zero
set $C$: this calculation takes place on $W$.

The two ordinary opens
\[
 \begin{split}
 W_{\rm pas}&=\{\psi<\Phi_0+\kappa\},\\
 W_{\rm two}&=\{\psi>\Phi_0+|a|+|b|\}
 \end{split}
 \tag{CS27}
\]
cover $W$ because $\varepsilon>0$, and the second contains every
passive-zero point.  In terms of continuous norm ratios they are
strict inequalities, hence are ordinary opens rather than rational
domains being substituted for actual sheaf opens.

On $W_{\rm pas}$, the cap equals $\psi$.  Equation (CS21) makes
both $\sum_eS_{e,s}$ and $\mathfrak a(a,b;R_*^{s-1})$ zero.
On $W_{\rm two}$ all passive entries are strictly inactive for
all four norms.  Indeed unit scaling changes a norm potential
by at most the absolute value of its unit valuation, so
\[
                 \Phi_i\le\Phi_0+|a|+|b|<\psi
                              \qquad(i=0,A,B,T).
 \tag{CS28}
\]
This holds at every analytic point over its completed residue
field.  The two-entry recurrence (CS18) therefore applies on
that entire open for $R_*$, including every cap-transition
curvature term and mixed derivative.  Actual sheaf locality gives
\[
                   \sum_eS_{e,s}=-\mathfrak a(a,b;R_*^{s-1})
                                      \quad\hbox{on all of }W.
 \tag{CS29}
\]
No triple has been extended from one member of (CS27): the cap
and every ensuing triple were already actual on all of $W$.
The cover only verifies equality between global actual formulas.

\subsection{The explicit whole-overlap triple}

For each oriented edge use $(p_i,p_j,c_e,R_*)$ in (CS6)--(CS9).
Put
\[
 Q_{\rm edge}=\sum_e\left\{(0,0,M_{e,s})+\frac12
     \bigl(Q_{e,s}^{\mathrm{tet}}-R_*Q_{e,s-1}^{\mathrm{tet}}\bigr)
                                                       \right\}.
 \tag{CS30}
\]
Then (CS10) and (CS29) give
\[
 P_s^*=\partial_2Q_{\rm edge}
                 +\frac12\mathfrak a(a,b;R_*^{s-1}).
 \tag{CS31}
\]
The alternating symbol and the extending pair differ by
\[
 \frac12\mathfrak a(a,b;R_*^{s-1})
       =e_{a,b}R_*^{s-1}
                   -\partial_2(0,0,abR_*^{s-1}/2).
 \tag{CS32}
\]
For $d=s-1\ge1$, set
\[
                h_*=-\Psi\sum_{\ell=0}^{d-1}R_*^\ell R^{d-1-\ell}.
 \tag{CS33}
\]
All curvatures are closed and $R_*-R=-D\Psi$, so
$Dh_*=R_*^d-R^d$.  The independently closed pair $e_{a,b}$
satisfies the exact general-degree CPT identity
\[
                 e_{a,b}R_*^d-e_{a,b}R^d
                                      =\partial_2H_{h_*}(e_{a,b}),
 \tag{CS34}
\]
where, explicitly,
\[
 H_{h_*}(e_{a,b})=
 \bigl(d''h_*\wedge(e_{a,b})',
       -d'h_*\wedge(e_{a,b})'',
                        h_*\wedge d'(e_{a,b})'\bigr).
 \tag{CS35}
\]
Here $h_*$ has degree $(d-1,d-1)$.  CPT is not being applied
to one of the nonclosed residual summands.

Define on the whole original $W$
\[
 \boxed{\quad
 q_s=Q_{\rm edge}-Q_{\rm met}
           -(0,0,abR_*^{s-1}/2)+H_{h_*}(e_{a,b}).
 \quad}
 \tag{CS36}
\]
Every term is an explicitly specified actual form.  Equations
(CS26), (CS31), (CS32), (CS34) prove (CS3), with these signs.
Its last pair extends through the common zero set because the
original $R,a,b$ are defined on all of $N$.

\begin{theorem}[All-rank actual relative commuting-square witness]
For every $r\ge2$, over the original field and absolute value,
the prescribed SHT square of two distinct-entry analytic unit
changes has the explicit actual witness (CS3).  For $r=2$
use frozen CDP2--CDP3.  For $r\ge3$ use (CS24)--(CS36).
All valuation walls, the arbitrary common metric, all mixed
base derivatives and zeros of proper subsets of the tuple
entries are retained.
\end{theorem}

No boundedness, fixed valuation order, rational-point test,
product family or analytic root is assumed.  Divisions by
$2,3,6$ are real coefficient operations and impose no residual
characteristic restriction.  The rank-three passive target is
$\mathbf P^0$; the proof handles that endpoint.  The rank-two
theorem is inherited, not counted again.

\subsection{Finite path orders and the original-support bridge}

For two paths consisting of the same prescribed diagonal unit
factors, an adjacent swap on distinct entries leaves all common
prefix and suffix terms identical.  Apply (CS3) to the actual
intermediate tuple.  A supplied finite sequence of such swaps
gives the literal identity
\[
                 p_1-p_0=\partial_2q+e|_W,\qquad
 q=\sum_{\rm swaps}\pm q_s,\qquad
 e=R^{r-2}\sum_{\rm swaps}\pm e_{a,b}.
 \tag{CS37}
\]
This covers reorderings of one factor per entry and shuffles
preserving the within-entry order of repeated factors.  It
does not assert same-entry merger relations or arbitrary
elementary-shear and pivot coherence.

For the actual Green consequence, retain original global regular
CI tuples $s_0,s_1$ of the same metrized line bundle on smooth
projective $X/K$.  Write
\[
 \begin{gathered}
 Z=[Y_1]-[Y_0],\quad T=|Z|,\quad S=|Y_0|\cup|Y_1|,\\
 U=X^{\rm an}\setminus T,\quad V=X^{\rm an}\setminus S,\quad
 C_U=S\cap U.
 \end{gathered}
 \tag{CS38}
\]
Suppose $N\subset U$ is an ordinary neighborhood of $C_U$ and
$N=N_0\cup N_1$.  Supply the local finite analytic frame paths
of frozen SHT/RSP, yielding
\[
       g|_{N_i\cap V}=\partial_1p_i+E_i,\qquad
       g=G_{s_1}-G_{s_0},\qquad E_i\in B^{s,s}(N_i).
 \tag{CS39}
\]
Assume their complete primitive expressions on the whole overlap
differ by a supplied finite sequence of the swaps just treated,
with all unit factors and intermediate tuples defined on the
unpunctured overlap.  Equation (CS37) supplies the actual triple
$q=(A_{01},B_{01},M_{01})$ and extending pair $e=(e',e'')$.

For a subordinate ordinary smooth partition $\chi_0+\chi_1=1$,
write $\chi=\chi_1$ and use actual zero extensions of all
overlap-supported products in
\[
 \begin{split}
 \alpha={}&\chi_0\alpha_0+\chi_1\alpha_1
                  +d'\chi\wedge A_{01}+d''\chi\wedge M_{01},\\
 \beta={}&\chi_0\beta_0+\chi_1\beta_1
                  +d''\chi\wedge B_{01}+d'\chi\wedge M_{01},\\
 E={}&\chi_0E_0+\chi_1E_1
                             -d'\chi\wedge e'-d''\chi\wedge e''.
 \end{split}
 \tag{CS40}
\]
The exact RAG6--RAG9 calculation gives
$g=d'\alpha+d''\beta+E$ on $N\cap V$, with $E$ actual on $N$.
In particular its differentiated extending-pair terms remain,
as do both mixed $D\chi\wedge M_{01}$ terms before their
opposite-sign cancellation.

Choose the $U$-only RSP2 cutoff $\zeta$, equal to one near
$C_U$ with closed $U$-support inside $N$.  Let
$\widetilde\alpha,\widetilde\beta$ be the actual zero extensions
of $\zeta\alpha,\zeta\beta$ to $V$.  Set
\[
 g_V=g-d'\widetilde\alpha-d''\widetilde\beta,\qquad
 g_Z=\text{the actual gluing of $g_V$ with $E$ near $C_U$}.
 \tag{CS41}
\]
This is an actual $B^{r-1,r-1}$ form on exactly $U$.
For every actual $\eta\in B_c^{n-r,n-r}(X^{\rm an})$, frozen
support avoidance gives $\operatorname{supp}D\eta\Subset V$.
The products of the correction primitives with $D\eta$ have
compact support in $V$, so full actual Stokes gives zero for
their integrals.  The original frozen CI8 identity yields
\[
                       \int_Ug_Z\wedge D\eta=-\int_Z\eta
                              \quad\hbox{for every such actual test}.
 \tag{CS42}
\]
No separated-test restriction, new growth condition, enlarged
deleted support or current-to-form inference is involved.

For more than two opens, RAG's sequential whole-overlap witness
must still be supplied after each partition gluing step.
Pairwise swaps are not asserted to supply all higher witnesses
for an arbitrary cover.  A global diagonal comparison already
has an old single-path repair.  Reordering that global path is
therefore not advertised as a new previously unsolved cycle
class.  The new assertion is the all-rank relative identity for
the formerly unresolved prescribed square and its two-open
actual repair interface.

\subsection{Actual refinements and remaining scope}

The object is a specified finite actual path pair; its equivalence
is witnessed by the actual triple (CS36) and the extending pair
(CS3).  No quotient or completion is used as a representative.
Restriction, finite sums, full differentiation and wedge products
are the operations.  Local presentations use the entire original
norm and unit data, the single cap minimum, and finite PPC/CPT/MT
formulas.  Refinement preserves every identity because scalar and
first-jet carrier restrictions were checked on the complete
refinement.

The pair and triple spaces retain their own actual coefficient-final
topologies.  For a fixed whole presentation the cap introduces one
finite arrangement in the stated unit coordinates with all real
intercepts; the remaining operations are finite polynomial and
differential operations with all mixed seminorms.  This gives
fixed-presentation finite-stage control only.  It does not identify
image-final and ambient inherited topologies or prove new global
strictness, Hausdorffness or bornologicity.

Nonproduct refinements such as $(T_1,T_2,T_1+T_2)$ are retained:
no independent-coordinate image replaces their actual relations.
Passive zeros of such expressions are explicitly handled by the
cap.  In moving-puncture curve families, all units and norm
potentials may depend on every base coordinate.  Their complete
derivatives remain, and any number of passive base directions may
be adjoined.  No fiberwise compact angular homotopy is used.

Every inequality and active-set statement holds at arbitrary
analytic points over their completed residue fields.  The
original $K$ and absolute value remain fixed.  No degree
division or higher actual-form descent is used.  ASP is respected
by the nonzero extending pair; LTF is respected because no
truncated pair is assumed closed.  The separate TROP-NR, CI0,
FL, JR, AQ, QP, ND and wedge failure scopes are unchanged.

General homological triviality still has no all-actual period
theorem here.  Arbitrary unequal-ideal cancellation, arbitrary
analytic matrix coherence, full lower-middle OPEN-KERNEL,
arbitrary higher actual descent and complete unconditional
Beilinson--Bloch height theory remain separate open tasks.
The original-support bridge above covers only the specified
original CI data with its finite overlap certificate.

\paragraph{Construction provenance.}
The diagonal constructor supplied the capped-passive proof, using
the audit/challenge constructor's normalized-edge derivation and
independent coefficient sign checks.  Both are construction
contributors and are excluded from independent reviews of this
module.  Root reviews this module without contributing a
construction step; MT is a separately audited dependency.
No external peer-review, formal-verification or novelty claim
is made.

